% Figure 21.2 : où le mot de passe du Secret colis-db se lit en clair (relevés du chapitre 21).
\documentclass[tikz,border=4pt]{standalone}
\usepackage{quai}
\begin{document}
\begin{tikzpicture}[x=1cm,y=1cm,
    e/.style={qbox, font=\scriptsize, text width=4.3cm, align=left, inner sep=4pt, minimum height=1.3cm},
    clair/.style={e, draw=QBrique, fill=QBriqueBg},
    b64/.style={e, draw=QOcre, fill=QOcreBg},
    prot/.style={e, draw=QVert, fill=QVertBg}]
  \node[b64] (a) at (0,1.6) {\textbf{{\ttfamily kubectl get secret -o yaml}}\\{\ttfamily Y2hhbmdlLW1vaS1l...}\\base64 : un simple codage, {\ttfamily base64 -d} le lit};
  \node[clair] (b) at (4.8,1.6) {\textbf{annotation}\\{\ttfamily last-applied-configuration}\\écrite par {\ttfamily kubectl apply} :\\le {\ttfamily stringData} en clair};
  \node[clair] (c) at (9.6,1.6) {\textbf{etcd}\\{\ttfamily /registry/secrets/ch21/...}\\stocké tel quel, sans chiffrement\\(chapitre 46)};
  \node[clair] (d) at (0,-0.3) {\textbf{dans le Pod, variable}\\{\ttfamily \$POSTGRES\_PASSWORD}\\visible par tout processus du conteneur};
  \node[prot] (f) at (4.8,-0.3) {\textbf{dans le Pod, fichier}\\{\ttfamily /run/secrets/colis-db/}\\tmpfs en mémoire, mode {\ttfamily 0400}};
  \node[clair] (g) at (9.6,-0.3) {\textbf{le manifeste}\\{\ttfamily secret-db.yaml} dans Git :\\à ne jamais faire\\(chiffrer, chapitre 46)};
\end{tikzpicture}
\end{document}
